\documentclass[a4paper,12pt]{article}
\usepackage[utf8]{inputenc}
\usepackage[T1]{fontenc}
\usepackage[ngerman]{babel}
\usepackage{amsmath}
\usepackage{amssymb}
\usepackage{enumitem}
\usepackage{geometry}
\geometry{a4paper, margin=2.5cm}
\usepackage{parskip}

\title{Einsendeaufgaben -- Aufgabe 3.1\\[0.5em]
\large 61211 Analysis}
\author{}
\date{}

\begin{document}
\maketitle

\section*{Aufgabe 1}

\begin{enumerate}[label=\alph*)]

    \item
    Wir zeigen die Stetigkeit von $f$ in $(0,0)$ mit Hilfe des $\varepsilon$-$\delta$-Kriteriums.

    Sei $\varepsilon > 0$ beliebig. Wir wählen $\delta = \min(1, 2\varepsilon)$.
    Sei $\binom{x}{y} \in \mathbb{R}^2 \setminus \{0\}$ mit $\left\| \binom{x}{y} \right\|_\infty < \delta$.

    \textit{Fall 1: $x^2 \leq y^2$}

    Wir zeigen zunächst, dass $\left| \frac{\sin(x)}{x} \right| \leq 1$ gilt.

    Wir betrachten $g(x) = x - \sin(x)$. $g'$ ist gegeben durch $g'(x) = 1 - \cos(x)$. Da $1 \geq \cos(x)$ (MG 13.1.3) folgt $g'(x) = 1 - \cos(x) \geq 0$. Nach 3.3.13 ist $g$ wegen $g'(x) \geq 0$ monoton wachsend. Wegen $g(0) = 0$ und der Monotonie von $g$ gilt demnach:
    \[
        g(x) = x - \sin(x) \geq 0 \quad \forall x \in [0, \infty)
    \]
    Daraus folgt
    \[
        x \geq \sin(x) \quad \forall x \in [0, \infty).
    \]
    Für $x < 0$ gilt $x \leq \sin(x)$, denn
    \[
        -y \leq -\sin(y) = \sin(-y) \quad \text{für alle } y \in [0, \infty)
    \]
    wegen $y \geq \sin(y)$.

    Aus $x \geq \sin(x)$ für alle $x \geq 0$ und $x \leq \sin(x)$ für alle $x < 0$ folgt
    \[
        \left| \frac{\sin(x)}{x} \right| \leq 1 \quad \forall x \in \mathbb{R} \setminus \{0\}.
    \]

    Damit können wir nun $|f(x,y)| < \varepsilon$ zeigen:
    \begin{align*}
        |f(x,y)| &= \left| \frac{\sin(x)^4 \, y^6}{x^2 + y^2} \right| \\
        &\leq \left| \frac{\sin(x)^4 \, y^6}{2 x^2} \right| && (\text{wegen } x^2 \leq y^2) \\
        &\leq \left| \frac{\sin(x)^2 \, y^6}{2} \right| && \left( \text{wegen } \left| \tfrac{\sin(x)}{x} \right| \leq 1 \right) \\
        &\leq \left| \frac{y^6}{2} \right| && (\text{wegen } |\sin(x)| \leq 1) \\
        &\leq \frac{|y|}{2} && (\text{wegen } |y| < \delta \leq 1) \\
        &< \varepsilon && (\text{wegen } |y| < \delta \leq 2\varepsilon)
    \end{align*}

    \textit{Fall 2: $y^2 < x^2$}

    Wir können direkt schließen:
    \begin{align*}
        |f(x,y)| &= \left| \frac{\sin(x)^4 \, y^6}{x^2 + y^2} \right| \\
        &\leq \left| \frac{\sin(x)^4 \, y^6}{2 y^2} \right| && (\text{wegen } y^2 < x^2) \\
        &= \left| \frac{\sin(x)^4 \, y^4}{2} \right| \\
        &\leq \left| \frac{y^4}{2} \right| && (\text{wegen } |\sin(x)| \leq 1) \\
        &< \frac{|y|}{2} && (\text{wegen } |y| < \delta \leq 1) \\
        &< \varepsilon && (\text{wegen } |y| < \delta \leq 2\varepsilon)
    \end{align*}

    \item
    \[
        D_1 f(x,y) = \frac{4 \cos(x) \sin(x)^3 \, y^6 \, (x^2 + y^2) - \sin(x)^4 \, y^6 \cdot 2x}{(x^2 + y^2)^2}
    \]
    \[
        D_2 f(x,y) = \frac{6 \sin(x)^4 \, y^5 \, (x^2 + y^2) - \sin(x)^4 \, y^6 \cdot 2y}{(x^2 + y^2)^2}
    \]

    Wir setzen $D_1 f(0,0) = 0$. Für $\binom{a}{b}$ mit $b \neq 0$ gilt $D_1 f(a,b) = (f \circ \iota_1)'(a)$. Für $b = 0$ gilt $f(a, 0) = 0$, also
    \[
        D_1 f(a, 0) = 0 = (f \circ \iota_1)'(a).
    \]
    Also ist $D_1 f(x,y)$ die Ableitung von $(f \circ \iota_1)(x)$.

    Nach der gleichen Argumentation ist $D_2 f(x,y)$ mit $D_2 f(0,0) = 0$ die Ableitung von $(f \circ \iota_2)(y)$.

    $f$ ist also in $(0,0)$ partiell differenzierbar.

    \item
    Wir zeigen die Stetigkeit der partiellen Ableitungen $D_1 f$ und $D_2 f$ in $(0,0)$. Nach Satz 3.6.4 ist $f$ dann total differenzierbar in $(0,0)$.

    Wir zeigen dies zunächst für $D_1 f$.

    Nach der gleichen Argumentation von a) gilt
    \[
        \lim_{(x,y) \to 0} \frac{\sin^3(x) \, y^6}{x^2 + y^2} = 0
    \]
    und
    \[
        \lim_{(x,y) \to 0} \frac{\sin^2(x) \, y^3}{x^2 + y^2} = 0.
    \]

    Daraus schließen wir
    \begin{align*}
        &\lim_{(x,y) \to 0} D_1 f(x,y) \\
        &= \lim_{(x,y) \to 0} 4 \cos(x) \cdot \lim_{(x,y) \to 0} \frac{\sin^3(x) \, y^6}{x^2 + y^2} \\
        &\quad + \lim_{(x,y) \to 0} 2x \cdot \left( \lim_{(x,y) \to 0} \frac{\sin^2(x) \, y^3}{x^2 + y^2} \right)^2 = 0.
    \end{align*}

    Nun zeigen wir noch Stetigkeit in $(0,0)$ für $D_2 f$.

    Nach der gleichen Argumentation von a) gilt
    \[
        \lim_{(x,y) \to 0} \frac{\sin^4(x) \, y^5}{x^2 + y^2} = 0
    \]
    und
    \[
        \lim_{(x,y) \to 0} \frac{\sin(x)^2 \, y^3}{x^2 + y^2} = 0.
    \]

    Daraus schließen wir:
    \begin{align*}
        &\lim_{(x,y) \to 0} D_2 f(x,y) \\
        &= 6 \cdot \lim_{(x,y) \to 0} \frac{\sin(x)^4 \, y^5}{x^2 + y^2} + \lim_{(x,y) \to 0} 2y \cdot \left( \lim_{(x,y) \to 0} \frac{\sin(x)^2 \, y^3}{x^2 + y^2} \right)^2 \\
        &= 0.
    \end{align*}

    Die partiellen Ableitungen sind also stetig in $(0,0)$ und nach Satz 3.6.4 ist $f$ demnach total differenzierbar.

\end{enumerate}

\end{document}
